\section{Varför ännu en söm}
\label{sec:whyseam}

\code{Interface} gjorde \code{Stub} och \code{Spi} utbytbara, och därmed \code{EchoNode} testbar.
Men \code{Spi} själv ligger \emph{under} den sömmen. Utan något mer går den bara att köra på riktig
hårdvara.

Räkna på vad det skulle betyda. \code{Spi} blir klar i \kapref{ch:driver}. Hårdvaran finns i
\kapref{ch:bringup}. Däremellan skulle du ha en oprövad klass som ingen kört en enda gång --- och
första gången den körs görs det på ett kort som samtidigt felsöks av en annan grupp, över fyra
ledningar som kanske sitter rätt.

Det är inte ett bra sätt att hitta ett bitskiftningsfel.

\begin{cppcode}
namespace driver::can
{
/**
 * @brief Moves single bytes to and from the CAN controller, and frames transactions.
 *
 *        Knows nothing about registers, command bytes or CAN: it only carries bytes.
 */
class ByteTransport
{
public:
    virtual ~ByteTransport() noexcept = default;

    virtual void select() noexcept   = 0;   /**< Assert SS; begin a transaction. */
    virtual void deselect() noexcept = 0;   /**< Release SS; end the transaction. */

    [[nodiscard]] virtual std::uint8_t transfer(std::uint8_t byte) noexcept = 0;
};
} // namespace driver::can
\end{cppcode}

Tre metoder. Det är hela lagret.

\textbf{Varför \code{transfer()} returnerar en byte.} SPI är full duplex; en signatur som vore
\code{void transfer(std::uint8_t)} skulle ljuga om vad hårdvaran gör.

\textbf{Varför \code{select()}/\code{deselect()} är egna metoder.} \code{SS}-ramningen är en del av
\emph{protokollet}, inte av hur en byte överförs. Transporten ska inte behöva veta att en
transaktion är fem bytes lång. Håll den dum, så gör den bara en sak --- och är lätt att
implementera rätt både på målet och i ett test.

\textbf{Varför \code{noexcept} överallt.} Lagret körs på ett friställt mål utan undantag. Att
skriva det i signaturen redan nu gör det omöjligt att smyga in något som kastar.

\section{Konstanterna}
\label{sec:constants}

\begin{cppcode}
class Spi final : public Interface
{
    // ...
private:
    /** Register indices, as sent in the command byte. */
    static constexpr std::uint8_t RegStatus{0U};
    static constexpr std::uint8_t RegTxId{1U};
    static constexpr std::uint8_t RegTxDlc{2U};
    static constexpr std::uint8_t RegTxDataLo{3U};
    static constexpr std::uint8_t RegTxDataHi{4U};
    static constexpr std::uint8_t RegTxSend{5U};
    static constexpr std::uint8_t RegRxId{6U};
    static constexpr std::uint8_t RegRxDlc{7U};
    static constexpr std::uint8_t RegRxDataLo{8U};
    static constexpr std::uint8_t RegRxDataHi{9U};
    static constexpr std::uint8_t RegRxAck{10U};
    static constexpr std::uint8_t RegErrorFlags{11U};
    static constexpr std::uint8_t RegTxAbort{12U};

    /** STATUS register bit masks. */
    static constexpr std::uint32_t StatusTxReady{1U << 0U};
    static constexpr std::uint32_t StatusRxValid{1U << 1U};
    static constexpr std::uint32_t StatusError{1U << 2U};

    /** Command byte: bit 7 set means write. */
    static constexpr std::uint8_t CmdWrite{0x80U};

    /** Value written to the write-triggered registers. */
    static constexpr std::uint32_t Trigger{0x1U};

    ByteTransport& myTransport; /**< Transport used for every register access. Not owned. */
};
\end{cppcode}

\textbf{Index, inte offset.} Registerkartans offsetkolumn är ett arv från den minnesmappade
varianten. Det som går ut på ledningen är indexet, alltså $\text{offset}/4$. Skriv konstanterna som
index direkt, så slipper du en division vid varje anrop och en klass av fel där någon dividerar
två gånger eller ingen gång.

\textbf{\code{1U << 0U} och inte \code{0x01U}.} Masken skrivs som en skiftning därför att
registerkartan är skriven i bitpositioner. \code{StatusRxValid{1U << 1U}} går att jämföra rad för
rad mot tabellen; \code{0x02U} gör det inte.

\section{De två hjälpfunktionerna}
\label{sec:helpers}

\begin{cppcode}
void Spi::writeReg(const std::uint8_t index, const std::uint32_t value) const noexcept
{
    myTransport.select();

    static_cast<void>(myTransport.transfer(static_cast<std::uint8_t>(CmdWrite | index)));
    static_cast<void>(myTransport.transfer(static_cast<std::uint8_t>(value >> 24U)));
    static_cast<void>(myTransport.transfer(static_cast<std::uint8_t>(value >> 16U)));
    static_cast<void>(myTransport.transfer(static_cast<std::uint8_t>(value >> 8U)));
    static_cast<void>(myTransport.transfer(static_cast<std::uint8_t>(value)));

    myTransport.deselect();
}

std::uint32_t Spi::readReg(const std::uint8_t index) const noexcept
{
    myTransport.select();

    static_cast<void>(myTransport.transfer(index));

    std::uint32_t value{};
    value |= static_cast<std::uint32_t>(myTransport.transfer(0U)) << 24U;
    value |= static_cast<std::uint32_t>(myTransport.transfer(0U)) << 16U;
    value |= static_cast<std::uint32_t>(myTransport.transfer(0U)) << 8U;
    value |= static_cast<std::uint32_t>(myTransport.transfer(0U));

    myTransport.deselect();
    return value;
}
\end{cppcode}

Läs \code{writeReg()} mot \code{81 00 00 01 23} i \kapref{ch:spi}. Fyra saker att lägga märke till:

\textbf{Varje \code{static_cast<void>} är avsiktlig.} \code{transfer()} är \code{[[nodiscard]]},
och under en skrivning är returvärdet meningslöst. Casten säger till både kompilatorn och läsaren
att det är känt.

\textbf{I \code{readReg()} kastas den första byten.} Slaven laddar sin svarsskiftare i \emph{slutet}
av kommandobyten. Fyra \code{transfer()} bidrar till värdet, inte fem.

\textbf{Kommandobyten är bara \code{index} vid läsning.} Bit 7 låg betyder läsning; ingen
\code{CmdWrite}.

\textbf{Metoderna är \code{const}.} De ändrar inte \code{Spi}-objektets tillstånd --- referensen i
sig ändras inte. Att hårdvaran ändras är en annan sak; \code{const} i C++ talar om objektet, inte
om världen. Det är därför \code{hasError() const} kan anropa \code{readReg()}.

\subsection{Teckenutvidgningsbuggen}

Det här är projektets vanligaste bugg, och den är värd en egen rubrik:

\begin{cppcode}
value |= myTransport.transfer(0U) << 24U;                             // fel
value |= static_cast<std::uint32_t>(myTransport.transfer(0U)) << 24U; // rätt
\end{cppcode}

\code{transfer()} returnerar en \code{std::uint8_t}. I den första raden befordras den till
\code{int} --- 32 bitar \emph{med tecken} --- och skiftas 24 steg. För ett värde med bit 7 satt
skriver det in i teckenbiten. Resultatet är odefinierat beteende, och i praktiken ett negativt tal
som sedan konverteras till \code{std::uint32_t} med \code{0xFF}-bytes där det inte ska finnas
några.

\textbf{Symptomet: drivrutinen fungerar för små värden och går sönder så fort en byte har bit 7
satt.} Alltså fungerar den för \code{0x123} och går sönder för \code{0xAABB}.

\begin{aside}
\note[Regel] Casta \textbf{före} skiftningen, inte efter. Samma regel gäller packningen i
\kapref{ch:driver}.
\end{aside}

\subsection{Varför hjälpfunktionerna är privata}

De är implementationsdetaljer. Ingen utanför klassen ska kunna peta i registren, av samma skäl som
\code{Interface} inte har en \code{readStatusRegister()}: den dag någon gör det i
applikationskoden är arkitekturen bruten.

Den utdelade testsviten testar dem ändå --- genom \code{send()} och \code{receive()}, och genom att
inspektera vilka bytes som kom ut. Att testa en privat metod direkt är nästan alltid ett tecken på
att den borde vara publik, eller på att testet testar fel sak.

\subsection{Vad lagret inte gör}

\code{readReg()} och \code{writeReg()} är rena översättningar av protokollet och innehåller ingen
logik alls: ingen validering, ingen statuskontroll, inget begrepp om vad ett register betyder.

All betydelse ligger i lagret ovanför, och det är \kapref{ch:driver}. Så fort \code{writeReg()}
börjar innehålla "om index är \code{RegTxSend}, kontrollera först att..." har två lager smält
ihop, och båda blev svårare att testa.
